\documentclass[11pt]{article}
\usepackage{icmlko}
\begin{document}

\title{바로 세운 넷은 잃고 안 건드린 둘이 벌었다 --- 그리고 씨앗 구간이 좁았던 건 팔이 짝지어져 있어서였다}
\author{월드모델 실험실}
\date{노트 609 · 610 · 2026년 8월 5일}
\maketitle

\begin{abstract}
노트 $608$ 이 ``도메인 수준에서 표본 SE 가 씨앗 SE 의 $12.6$ 배''를 냈고,
남은 것에 \textbf{판 자체는 안 쟀다}고 적었다. 판($3{,}369$ 행)에서 재니
비가 팔마다 다르다.

\begin{center}
\begin{tabular}{lrrr}
\hline
팔 & 씨앗 SE & 표본 SE & 비 \\
\hline
위키 가림(\textbf{같은 적합}) & $0.00047$ & $0.00233$ & $\mathbf{5.0}$ \\
채택(\textbf{따로 적합}) & $\mathbf{0.00235}$ & $0.00326$ & $1.4$ \\
\hline
\end{tabular}
\end{center}

\textbf{씨앗 SE 가 다섯 배 다르다.} 가림 팔은 두 팔이 \emph{같은 모형}을
쓰므로 씨앗 잡음이 상쇄된다. \textbf{``씨앗 구간이 좁다''는 효과가
확실하다는 뜻이 아니라 팔이 짝지어져 있다는 뜻이다} --- 그리고
\textbf{채택은 언제나 따로 적합하는 팔이다.}
\end{abstract}

\begin{figure}[t]
\centering
\includegraphics[width=\textwidth]{figs/f1.pdf}
\caption{\textbf{가} --- 씨앗 폭은 팔 설계가 정한다.
\textbf{나} --- 빨강이 실제로 고친 넷이다.
\textbf{다} --- 고친 넷만 음수다.}
\end{figure}

\section{이번 세션의 채택을 그 자로 다시 쟀다}

노트 $583/586$ 이 \texttt{axis\_orient.json} 을 모형 경로에 물려
\texttt{entry\_friction} 을 도서 $\cdot$ 모바일 $\cdot$ 애니 $\cdot$ 펀딩에서
원래 방향으로 되돌렸다. \textbf{판 $+0.0019$}(씨앗 $40$ 중 $34$)로 채택했다.

\begin{center}
\begin{tabular}{lr}
\hline
 & 판 차 \\
\hline
평균(씨앗 $8$) & $+0.00182$ \\
\textbf{씨앗} $95\%$ & $[\mathbf{-0.00166}, +0.00443]$ \\
\textbf{표본} $95\%$ & $[\mathbf{-0.00459}, +0.00801]$ \\
\hline
\end{tabular}
\end{center}

\textbf{씨앗 기준으로도 이미 $0$ 을 문다.} 씨앗 $40$ 의 \emph{평균}은 확실히
양수지만($34/40$ 은 우연이 $10^{-5}$), 씨앗을 늘려도 \textbf{표본 성분은
안 줄어든다} --- 순수 표본 SD 가
$\sqrt{0.00326^2-0.00235^2}=\mathbf{0.00226}$ 이라 효과 대비 $t{=}0.8$ 이다.

\begin{quote}
\textbf{채택은 유지한다.} 되돌릴 근거는 ``판이 내려간다''이지 ``구간이
넓다''가 아니고(노트 $524$), 집 밖 두 짝이 $40/40$ 으로 크다(KR 만화
$0.638\!\to\!0.684$ $\cdot$ 비게임 앱 $0.290\!\to\!0.505$). \textbf{바뀌는
것은 판 이득을 단정하지 않는다는 것}이다.
\end{quote}

\section{그런데 이득이 어디서 오나}

\begin{center}
\begin{tabular}{lrl}
\hline
 & 판 기여 & 표본 $95\%$ \\
\hline
\textbf{실제로 고친 넷} & $\mathbf{-0.00094}$ & $[-0.0047, +0.0026]$ \\
팝업($65$ 행) & $+0.00156$ & $[-0.0008, +0.0042]$ \\
웹툰($650$ 행) & $+0.00155$ & $[-0.0008, +0.0038]$ \\
안 고친 일곱 & $+0.00279$ & $[-0.0019, +0.0071]$ \\
\hline
\end{tabular}
\end{center}

\textbf{축을 바로 세운 네 도메인은 합쳐서 잃는다.} 판 이득 전부가
\emph{그 축을 건드리지도 않은} 도메인에서 온다 --- 팝업과 웹툰 둘이
$+0.0031$ 이고, 그 둘의 \texttt{entry\_friction} 은 방향이 바뀌지 않았다
(팝업은 애초에 없고 웹툰은 노트 $262$ 가 막았다).

\begin{quote}
\textbf{이웃 효과다}(노트 $551 \cdot 525 \cdot 559 \cdot 560 \cdot 563
\cdot 596 \cdot$ 여기 --- \textbf{일곱 번째}). 한 도메인의 축을 고치면
\emph{공유 모형}이 바뀌고 그것이 다른 도메인 예보에 실린다. 이번에는
\textbf{그 몫이 직접 몫보다 크고 부호도 반대}다.
\end{quote}

\section{채택의 근거가 틀렸다는 뜻은 아니다}

논문 $452$ 가 채택을 ``두 해 동안 거꾸로 꽂혀 있던 축을 바로 세웠다''로
냈다. \textbf{진단은 여전히 옳다} --- \texttt{axis\_orient.json} 이
$2024$ 년부터 기록해 둔 뒤집힘이 모형 경로에서 안 읽히고 있었고, 그것이
훈련 가중 상관을 $-0.023$(방향 없음)으로 만들어 \texttt{OOC\_DROP} 이
존재할 이유가 됐다. 되돌리자 그 가림이 필요 없어졌고 \textbf{집 밖 두 짝이
크게 올랐다.}

\begin{quote}
\textbf{바뀐 것은 ``판이 왜 올랐나''다.} 논문 $452$ 는 판 $+0.0019$ 를
\emph{고친 도메인들이 제대로 예보하게 됐다}로 읽었다. 실제로는 고친 넷이
$-0.0009$ 이고 이웃이 $+0.0031$ 이다. \textbf{옳은 진단이 옳은 인과를
보장하지 않는다.}
\end{quote}

\section*{조항}

\begin{quote}
\textbf{씨앗 구간의 폭은 효과의 확실함이 아니라 팔의 짝지어짐을 잰다.}
같은 적합에서 열만 가리는 팔은 씨앗 SE 가 $0.00047$ 이고, 두 번 적합하는
팔은 $0.00235$ 다 --- \textbf{다섯 배}. 이 판이 ``$12/12$'' $\cdot$
``$0$ 밖''으로 읽어 온 것 중 \emph{가림 팔}의 좁은 구간은 \textbf{짝지어짐
덕}이었다. \textbf{채택은 언제나 재적합 팔이므로 그 혜택을 못 받는다.}
\end{quote}

\section*{남은 것}

\begin{itemize}
\item \textbf{공유 인쇄기가 없다.} A/B 표의 ``$95\%$ 구간''이 노트별
스크립트에 복제돼 있어 규약을 한 곳에서 못 바꾼다. \textbf{노트 $607$ 이
겪은 것과 같은 자리}다(조항을 한 함수에 넣으면 다른 함수는 모른다).
\item \textbf{팝업 $+0.0809$ 가 왜 그렇게 큰가} --- $65$ 행이고 그중
$7$ 행은 축이 아예 없다(노트 $602$). 이웃 효과가 \emph{빈 곳}에서 제일
크게 나타나는지 안 봤다.
\item \textbf{유의성 위 선택 축소 추정 $\cdot$ 군집 부트스트랩}이
노트 $608$ 에서 그대로 넘어왔다. 안 했다.
\end{itemize}

\end{document}
